\section*{Executive summary}
\addcontentsline{toc}{section}{Executive summary}
\label{sec:execsummary}

\paragraph{What this note is.}
This note documents the application of \emph{unbinned machine-learning
unfolding} (OmniFold) to MINERvA's medium-energy neutrino data --- first as a
faithful reproduction of a published flagship measurement, then as its
extension to simultaneous three-, four- and five-observable cross sections
that no binned method has produced at MINERvA. It is written so that the
first three pages give the complete story; everything afterwards is the
supporting technical record.

\paragraph{The method, in plain language.}
Traditional unfolding corrects detector effects with a migration matrix
between histogram bins, so each added observable multiplies the number of
bins and the method drowns in dimensions. OmniFold instead trains a
classifier to learn, event by event, a weight that morphs the simulation
into the data, and pulls that weight back to the true (pre-detector) level
through the simulation's truth--reco pairing. No binning is imposed until
the very end, so adding an observable is just adding an input column.
This analysis uses gradient-boosted decision trees as the learners and
validates the choice extensively (closure tests, coverage, calibrated-
classifier checks, neural-network and transformer cross-checks).

\paragraph{Headline results.}
\begin{itemize}[itemsep=2pt]
  \item \textbf{Reproduction (the demonstrator).} Our unbinned
        $d^{2}\sigma/(d\pT\,d\pPar)$ matches the published measurement
        (arXiv:2106.16210): total cross-section ratio $1.011$,
        \SI{94}{\percent} of bins within \SI{10}{\percent}, per-bin pulls
        sub-$1\sigma$, and a fully independent uncertainty budget whose
        median per-bin uncertainty (\SI{6.87}{\percent}) matches the
        published \SI{6.86}{\percent} (\S\ref{sec:results},
        \S\ref{sec:systematics}, Fig.~\ref{fig:xsec}).
  \item \textbf{First 3D/4D/5D MINERvA unfolds.} Adding available energy
        ($\Eavail$), momentum transfer ($q_{3}$), and hadronic mass ($W$)
        produces the first simultaneous three-, four- and five-observable
        unfolds of MINERvA data; every result reproduces its
        lower-dimensional predecessor when marginalised --- a built-in
        cross-check (\S\ref{sec:3d}).
  \item \textbf{Known physics recovered independently.} All four truth-level
        generators (GENIE, MINERvA Tune, NuWro, GiBUU) under-predict the
        data at low $\Eavail$; pion final-state-interaction dials move the
        spectrum by $<\SI{1}{\percent}$ while adding Valencia 2p2h fills the
        quasielastic--$\Delta$ dip --- the established MINERvA
        low-recoil/2p2h deficit, re-found by an independent method
        (\S\ref{sec:3d-2p2h}, Fig.~\ref{fig:3dmodels}).
  \item \textbf{A new localisation of the remaining excess.} The
        high-$\Eavail$ excess that 2p2h cannot explain is resolved in the
        $(\Eavail,W)$ plane: it concentrates in the deep-inelastic corner
        ($W\ge\SI{1.8}{GeV}$), where all four generators miss the data at
        $9$--$18\sigma$ on a dedicated covariance --- no current model
        reproduces it (\S\ref{sec:eavailw}, Fig.~\ref{fig:eavailWband}).
  \item \textbf{A methodological finding.} A 160-throw ``unified-throw''
        test of the standard add-the-blocks covariance construction shows
        it \emph{underestimates} the 4D systematic by a factor $\sim2$ in
        the strongest-migration corner; the published 4D covariance adopts
        the corrected magnitude (\S\ref{sec:uq-combine}).
  \item \textbf{Toward raw-data unfolding.} A point-cloud transformer
        unfolding the actual reconstructed calorimeter clusters now yields
        an absolutely-normalised, full-statistics cross section, unbiased in
        closure to $\sim\SI{1}{\percent}$ and agreeing with the production
        result at the \SI{9}{\percent} (training-limited) level
        (\S\ref{sec:pet}).
  \item \textbf{A first full-phase-space cross section.} Removing the muon
        kinematic cuts (which hide a third of the signal rate) gives
        $\sigma_{\mathrm{FPS}}=\SI{4.5e-38}{cm^2/nucleon}$, anchored to the
        restricted measurement at \SI{0.6}{\percent} where they overlap,
        with a three-prior envelope declaring the extrapolation cost
        cell-by-cell (\S\ref{sec:fps}; systematics in production).
\end{itemize}

\paragraph{How to read this note.}
Sections~\ref{sec:experiment}--\ref{sec:method} define the samples and the
method; \S\ref{sec:results} is the 2D reproduction;
\S\ref{sec:systematics} and App.~\ref{app:statmethods} are the uncertainty
machinery (the appendix is self-contained and tutorial-grade);
\S\ref{sec:validation} is the validation suite; the physics is in
\S\ref{sec:3d}--\ref{sec:eavailw}; \S\ref{sec:pet} is the point-cloud
track; \S\ref{sec:summary} states the planned full-phase-space extension.
Appendix~\ref{app:openq} records the resolution of all questions raised in
internal review. Every number quoted here is cross-checked in the
repository's validation ledger.
